\documentclass[11pt]{article}
\usepackage[margin=1in]{geometry}
\usepackage[shortlabels]{enumitem}
\usepackage{amsmath, amssymb, amsthm}
\usepackage{newpxtext, newpxmath}
\usepackage{hyperref, cleveref}

\newtheorem{theorem}{Theorem}[section]
\newtheorem{lemma}{Lemma}[section]
\newtheorem{corollary}{Corollary}[section]
\newtheorem{prop}{Proposition}[section]

\theoremstyle{definition}
\newtheorem{definition}{Definition}[section]
\newtheorem{example}{Example}[section]

\theoremstyle{remark}
\newtheorem{remark}{Remark}[section]

\title{\textbf{Introduction to Topology}}
\author{Alex Davison}

\begin{document}
\maketitle

\tableofcontents

\section{Introduction}

\emph{Topology} studies mathematical spaces where quantitative features such as distance, angle, volume, and area are ignored, and qualitative features such as openness, continuity, connectedness, compactness, and holes are focused on instead. Given a structureless set $A$, the minimum amount of data you can append to $A$ such that it properly becomes a \emph{space} with these most basic qualities is exactly a topology.

These notes introduce general topology and some topics of algebraic topology suitable for a strong undergraduate. Some exposure to topological topics from, e.g., an analysis course is helpful. We begin with metric spaces, where a notion of distance is still present; the open sets of a metric space are the motivating example of the abstract topologies defined afterwards.

\section{Metric spaces}

\begin{definition}[metric spaces]
    Given a set $X$, a function $d : X \times X \to [0,\infty)$ is a \emph{metric} if these hold for all $x, y, z \in X$:
    \begin{enumerate}[(a)]
        \item $d(x, y) = 0$ iff $x = y$ (identity of indiscernibles);
        \item $d(x, y) = d(y, x)$ (symmetry);
        \item $d(x, z) \leq d(x, y) + d(y, z)$ (triangle inequality).
    \end{enumerate}
    A set together with a metric $(X, d)$ is called a \emph{metric space}.
\end{definition}

\begin{example}
    $\mathbb{R}$ has the \emph{standard metric} $d(x, y) = |x - y|$, and more generally $\mathbb{R}^n$ has the Euclidean metric $d(x, y) = \sqrt{\sum_{i = 1}^n (x_i - y_i)^2}$. Any set $X$ also admits the \emph{discrete metric}
    \[
    d(x, y) =
    \begin{cases}
        0 & \text{if } x = y, \\
        1 & \text{if } x \neq y,
    \end{cases}
    \]
    in which any two distinct points are distance $1$ apart.
\end{example}

\begin{definition}[balls]
    Let $(X, d)$ be a metric space. The \emph{open ball} centered at a point $x_0 \in X$ with radius $r > 0$, denoted $B_r(x_0)$, is
    \[
    B_r(x_0) = \{x \in X \mid d(x_0, x) < r\}.
    \]
    Replacing $<$ by $\leq$ yields the \emph{closed ball} $\overline{B_r}(x_0)$.
\end{definition}

\begin{definition}[open sets in metric spaces]
    \label{metricopen}
    A set $U \subseteq X$ is \emph{open} in a metric space $(X, d)$ if for every $x \in U$ there is some $\varepsilon > 0$ with $B_\varepsilon(x) \subseteq U$.
\end{definition}

\begin{definition}[neighborhoods]
    A \emph{neighborhood} of a point $x$ in a metric space is an open set $U$ containing it, $x \in U$.
\end{definition}

\begin{remark}
    Some authors call any set that merely \emph{contains} an open set around $x$ a neighborhood of $x$; in these notes neighborhoods are always open.
\end{remark}

\begin{lemma}
    \label{ballsopen}
    Every open ball is an open set.
\end{lemma}

\begin{proof}
    Let $y \in B_r(x_0)$ and set $\varepsilon = r - d(x_0, y) > 0$. If $z \in B_\varepsilon(y)$, then the triangle inequality gives
    \[
    d(x_0, z) \leq d(x_0, y) + d(y, z) < d(x_0, y) + \varepsilon = r,
    \]
    so $B_\varepsilon(y) \subseteq B_r(x_0)$. Hence every point of $B_r(x_0)$ has a ball around it contained in $B_r(x_0)$, which is therefore open.
\end{proof}

\begin{prop}
    \label{unionballs}
    A set in a metric space is open iff it is a union of open balls.
\end{prop}

\begin{proof}
    If $U = \bigcup_\alpha B_\alpha$ is a union of open balls and $x \in U$, then $x \in B_\alpha$ for some $\alpha$, and since $B_\alpha$ is open (\cref{ballsopen}) some ball around $x$ lies in $B_\alpha \subseteq U$; hence $U$ is open. Conversely, if $U$ is open, then for every $x \in U$ choose a ball $B_x$ with $x \in B_x \subseteq U$; then $U = \bigcup_{x \in U} B_x$ is a union of open balls.
\end{proof}

\begin{example}
    The open balls of $(\mathbb{R}, |\cdot|)$ are exactly the open intervals, so by \cref{unionballs} a set is open in $\mathbb{R}$ iff it is a union of open intervals:
    \[
    \tau_{\mathbb{R}} = \left\{\bigcup_{\alpha \in A} (a_\alpha, b_\alpha) \;\middle|\; (a_\alpha, b_\alpha) \subseteq \mathbb{R} \text{ open intervals}\right\}.
    \]
    In the discrete metric every singleton $\{x\} = B_{1/2}(x)$ is an open ball, so every subset of $X$ is open.
\end{example}

\begin{definition}[continuity at a point]
    \label{localcontmet}
    A function $f : X \to Y$ between metric spaces $(X, d_X)$ and $(Y, d_Y)$ is \emph{continuous at} $x \in X$ if for every $\varepsilon > 0$ there is some $\delta > 0$ such that
    \[
    f(B_\delta(x)) \subseteq B_\varepsilon(f(x)),
    \]
    i.e., $d_X(x, x') < \delta$ implies $d_Y(f(x), f(x')) < \varepsilon$. The function $f$ is \emph{continuous} if it is continuous at every point of $X$.
\end{definition}

\begin{theorem}
    \label{pointwise}
    A function $f : X \to Y$ between metric spaces is continuous at $x \in X$ iff for every neighborhood $V$ of $f(x)$ there is a neighborhood $U$ of $x$ with $U \subseteq f^{-1}(V)$, equivalently with $f(U) \subseteq V$.
\end{theorem}

\begin{proof}
    Suppose $f$ is continuous at $x$, and let $V$ be a neighborhood of $f(x)$. Since $V$ is open and contains $f(x)$, there is some $\varepsilon > 0$ with $B_\varepsilon(f(x)) \subseteq V$. By \cref{localcontmet} there is some $\delta > 0$ with $f(B_\delta(x)) \subseteq B_\varepsilon(f(x)) \subseteq V$, so $U = B_\delta(x)$ is a neighborhood of $x$ with $U \subseteq f^{-1}(V)$.

    Conversely, suppose the neighborhood condition holds, and let $\varepsilon > 0$. Then $V = B_\varepsilon(f(x))$ is a neighborhood of $f(x)$, so there is a neighborhood $U$ of $x$ with $U \subseteq f^{-1}(B_\varepsilon(f(x)))$. Since $U$ is open and contains $x$, there is some $\delta > 0$ with $B_\delta(x) \subseteq U$, hence $f(B_\delta(x)) \subseteq B_\varepsilon(f(x))$.
\end{proof}

\begin{theorem}
    \label{contpreimage}
    A function $f : X \to Y$ between metric spaces is continuous iff $f^{-1}(V)$ is open in $X$ for every open $V \subseteq Y$.
\end{theorem}

\begin{proof}
    Suppose $f$ is continuous and $V \subseteq Y$ is open. For any $x \in f^{-1}(V)$, the set $V$ is a neighborhood of $f(x)$, so by \cref{pointwise} there is a neighborhood $U$ of $x$ with $U \subseteq f^{-1}(V)$; thus every point of $f^{-1}(V)$ has a ball around it contained in $f^{-1}(V)$, which is therefore open.

    Conversely, suppose preimages of open sets are open, and let $x \in X$. If $V$ is a neighborhood of $f(x)$, then $f^{-1}(V)$ is open by assumption and contains $x$, so it is itself a neighborhood of $x$ contained in $f^{-1}(V)$; by \cref{pointwise}, $f$ is continuous at $x$.
\end{proof}

\Cref{contpreimage} says that continuity between metric spaces can be detected using only their open sets, without any mention of distances. This motivates taking open sets themselves as the fundamental structure.

\section{Topological spaces}

\begin{definition}[topological spaces]
    \label{topspace}
    Given a set $X$, a \emph{topology} on $X$ is a collection $\tau \subseteq 2^X$ of subsets of $X$ such that:
    \begin{enumerate}[(1)]
        \item the empty set and $X$ are open, i.e., $\emptyset, X \in \tau$;
        \item arbitrary unions of open sets are open, i.e., $\bigcup_{\alpha} U_\alpha \in \tau$ for any family $\{U_\alpha\}_\alpha \subseteq \tau$;
        \item the intersection of two open sets is open, i.e., $U \cap V \in \tau$ for any $U, V \in \tau$.
    \end{enumerate}
    A set together with a topology $(X, \tau)$ is called a \emph{topological space}. The elements of $\tau$ are its \emph{open sets}.
\end{definition}

\begin{prop}[metric topology]
    \label{metrictop}
    The collection of open sets of a metric space $(X, d)$ is a topology on $X$, called the \emph{metric topology} induced by $d$.
\end{prop}

\begin{proof}
    The set $\emptyset$ is open vacuously, and $X$ is open since $B_1(x) \subseteq X$ for every $x \in X$. If $x \in \bigcup_\alpha U_\alpha$ for a family of open sets, then $x \in U_\alpha$ for some $\alpha$, so some ball around $x$ lies in $U_\alpha$ and hence in the union; arbitrary unions of open sets are therefore open. Finally, if $x \in U \cap V$ with $U, V$ open, then there are balls $B_\varepsilon(x) \subseteq U$ and $B_\delta(x) \subseteq V$, so $B_{\min(\varepsilon, \delta)}(x) \subseteq U \cap V$; the intersection of two open sets is therefore open.
\end{proof}

\begin{remark}
    Not every topology arises from a metric; a space whose topology is induced by some metric is called \emph{metrizable}.
\end{remark}

\begin{example}
    The \emph{discrete topology} on a set $X$ is given by $2^X$, i.e., every subset of $X$ is open; it is induced by the discrete metric. The \emph{trivial} (or \emph{indiscrete}) \emph{topology} has only $\emptyset$ and $X$ open.
\end{example}

\begin{corollary}
    \label{fininteropen}
    The intersection of finitely many open sets is open.
\end{corollary}

\begin{proof}
    Induct on axiom (3) of \cref{topspace}: the two-set case is the axiom itself, and if $U_1, \ldots, U_n$ are open, then $U_1 \cap \cdots \cap U_n = (U_1 \cap \cdots \cap U_{n-1}) \cap U_n$ is the intersection of two open sets.
\end{proof}

\begin{definition}[neighborhoods]
    As in the metric case, a \emph{neighborhood} of a point $x$ in a topological space is an open set containing it.
\end{definition}

\begin{definition}[closed sets]
    A set $F \subseteq X$ is \emph{closed} if the complement $X \smallsetminus F$ is open.
\end{definition}

\begin{corollary}
    Finite unions of closed sets are closed, and arbitrary intersections of closed sets are closed.
\end{corollary}

\begin{proof}
    If $F_1, \ldots, F_n$ are closed sets in $X$, then
    \[
    X \smallsetminus (F_1 \cup \cdots \cup F_n) = (X \smallsetminus F_1) \cap \cdots \cap (X \smallsetminus F_n)
    \]
    is open, being the intersection of finitely many open sets (\cref{fininteropen}), so $F_1 \cup \cdots \cup F_n$ is closed. If $\{F_\alpha\}_\alpha$ is any family of closed sets, then
    \[
    X \smallsetminus \bigcap_\alpha F_\alpha = \bigcup_\alpha (X \smallsetminus F_\alpha)
    \]
    is open, being the union of open sets, so $\bigcap_\alpha F_\alpha$ is closed.
\end{proof}

\end{document}
